% =============================================================================
\section{Experiments}
\label{sec:experiments}
\budget{3.0 pages, including Figures~\ref{fig:ondevice}, \ref{fig:samples} and the two tables}

\subsection{Setup}
\label{sec:setup}
\begin{plan}
  \item \textbf{Data.} MNIST \citep{lecun1998gradient}, Fashion-MNIST
        \citep{xiao2017fashion}, CelebA \citep{liu2015deep} as grey centre crops
        at $32\times32$ and $64\times64$. Pixels in $[0,1]$ with
        $\mathcal{N}(0,0.01^{2})$ noise, one global normalisation.
  \item \textbf{Model.} One architecture for all four: $M=\Mpatch$,
        \nSlices{} slices, $T=0.02$, $\tau=5$; coarse patches from $t=4.49$,
        fine ones from $t=2.22$. Two things differ. The patch sizes follow the
        resolution: $16\times16$ every 4 pixels and $8\times8$ every 2 pixels
        up to $32\times32$; $32\times32$ every 8 and $16\times16$ every 4 at
        $64\times64$. The global unit has $K=\Kglobal$ states for faces and
        $K=\KglobalClasses$ for MNIST and Fashion-MNIST; Section
        \ref{sec:quality} shows why.
        State the rule by which $K$ was chosen: a coarse search with the patch
        units fixed ($K\in\{16,128,784\}$ on MNIST and Fashion-MNIST,
        $K\in\{16,784\}$ on CelebA $32\times32$), taking the smallest $K$
        within one FID of the largest. Say that the choice was made on the
        reported metric with one seed, and give all values in
        Appendix~\ref{app:model-search}.
        \todo{CelebA $64\times64$ runs with $K=\Kglobal$.}
  \item \textbf{Training.} Each slice separately, from $t=\tau$ down, warm
        started from the previous slice; 512 fresh noised images per step.
        Three runs per data set: exact hidden averages; sampled with
        $S=\Slarge$; sampled with $S=\Ssmall$.
  \item \textbf{Sampling.} Noise start, reverse SDE \eqref{eq:reverse}, quench.
        \nSamples{} samples per model.
  \item \textbf{Metric.} Inception FID against the train split
        \citep{heusel2017gans}; the floor is the FID of real test images at the
        same sample count. \todo{say whether FID at n or the extrapolated
        FID$_\infty$ \citep{chong2020effectively} is reported.}
        State the protocol in full, because FIDs of different papers are not
        comparable without it: the original Inception graph, 8-bit images,
        grey replicated to three channels, the full train split as reference,
        the sample count.
  \item \textbf{What the FID does not show; say it.} On MNIST the Inception
        FID follows stroke texture and faint pixel noise more than digit
        identity: the quench alone moves it from about 50 to 14 at an
        unchanged share of clear digits, and faint noise on real test images
        gives 38. So show uncurated samples next to it, and keep the
        classifier-based numbers that exist (share of samples assigned to a
        class with confidence) as a supplementary column in the appendix.
        No further metrics are planned.
  \item \textbf{Compute.} One GPU per run; hours in
        Appendix~\ref{app:details}.
\end{plan}

\subsection{Training from Sampled Hidden States}
\begin{plan}
  \item \textbf{Main result} (Figure~\ref{fig:ondevice},
        Table~\ref{tab:results}; first wave, one seed, 2{,}000 samples): with
        $S=\Slarge$ the FID is 1 to 2 above exact training. CelebA
        $32\times32$: \FIDcelebLowSlarge{} against \FIDcelebLowExact; MNIST:
        \FIDmnistSlarge{} against \FIDmnistExact; Fashion-MNIST:
        \FIDfashionSlarge{} against \FIDfashionExact. With the corrector the
        gap on CelebA and MNIST is 0.3 and 0.6. Do not write ``matches'': the
        gap is several standard deviations (bootstrap 0.2 to 0.6).
        \todo{CelebA $64\times64$ when it is in.}
  \item \textbf{Fewer samples cost nothing.} $S=\Ssmall$ equals $S=\Slarge$ on
        all three data sets (CelebA \FIDcelebLowSsmall, MNIST \FIDmnistSsmall,
        Fashion \FIDfashionSsmall). So the gap to exact training does not come
        from the number of reads. \todo{say where it comes from: the earlier
        decomposition on a larger model put part of it on the visit-count rule
        in place of Adam; a run with fewer reads ($S=4$) would show where
        sampling starts to cost. Not run.}
  \item \textbf{Same samples.} With the same start and noise, exact and sampled
        training draw the same face, digit or garment in almost every column
        (the strip exact / $S=\Slarge$ / $S=\Ssmall$; Figure~\ref{fig:samples}
        or Appendix~\ref{app:more-results}). This is the strongest picture of
        the claim: use it.
  \item \textbf{State.} With sampled training on MNIST, 1 to 9 of 2{,}000
        trajectories diverged and were left out; none with exact training and
        none on the other data sets. \todo{cause not known.}
  \item \textbf{Seeds.} Spread over three seeds: \seedSdCelebLow{} (CelebA),
        \seedSdMnist{} (MNIST), \seedSdFashion{} (Fashion). \todo{second wave;
        the gap to exact training needs seeds of the exact runs as well.}
\end{plan}

% ---------------------------------------------------------------- Figure 2 ---
\begin{figure}[t]
  \centering
  \placeholderbox{\columnwidth}{2.1in}{%
    \textbf{Figure 2 placeholder.}\par\smallskip
    FID against the number of sampled hidden states.\par\smallskip
    x-axis: $S=\Ssmall$, $S=\Slarge$, exact ($S\to\infty$).\par
    y-axis: FID (log scale if the data sets differ much).\par
    One line per data set: MNIST, Fashion-MNIST, CelebA $32\times32$,
    CelebA $64\times64$; or four small panels.\par
    Dashed line: exact training. Error bars: seeds, where run.\par\smallskip
    Source: the lock-in runs; numbers as in Table~\ref{tab:results}.}
  \caption{\textbf{Training from sampled hidden states stays within 1 to 2 FID
    of training with exact hidden averages,} and 16 reads per clamped image do
    as well as 64. \todo{caption: what is plotted, the protocol in one line.
    With two values of $S$ that give the same FID the plot is two flat points
    per data set: consider the same-trajectory strip as this figure instead,
    or add a run with fewer reads.}}
  \label{fig:ondevice}
\end{figure}

% ----------------------------------------------------------------- Table 1 ---
\begin{table}[t]
  \caption{\textbf{FID on four data sets} (\nSamples{} samples, quench; one
    seed). Top: this work with exact hidden averages and with $S$ sampled
    hidden states; $K=\KglobalClasses$ for MNIST and Fashion-MNIST,
    $K=\Kglobal$ for CelebA. Bottom: real test images at the same sample
    count, the floor of this protocol.}
  \label{tab:results}
  \centering
  \small
  \setlength{\tabcolsep}{3.5pt}
  \begin{tabular}{lcccc}
    \toprule
      & MNIST & Fashion & CelebA\,32 & CelebA\,64 \\
    \midrule
    \multicolumn{5}{l}{\emph{This work}} \\
    exact averages        & \FIDmnistExact  & \FIDfashionExact  & \FIDcelebLowExact  & \FIDcelebHighExact  \\
    sampled, $S=\Slarge$  & \FIDmnistSlarge & \FIDfashionSlarge & \FIDcelebLowSlarge & \FIDcelebHighSlarge \\
    sampled, $S=\Ssmall$  & \FIDmnistSsmall & \FIDfashionSsmall & \FIDcelebLowSsmall & \FIDcelebHighSsmall \\
    $S=\Slarge$, corrector & \FIDmnistSlargeCorr & \FIDfashionSlargeCorr & \FIDcelebLowSlargeCorr & \FIDcelebHighSlargeCorr \\
    \midrule
    % FIDs reported in other papers are not in this table (2026-11-22): their
    % sample counts and implementations are not stated. See Section 5.3.
    real test images      & \floorMnist & \floorFashion & \floorCelebLow & \floorCelebHigh \\
    \bottomrule
  \end{tabular}
\end{table}

\subsection{Sample Quality and Comparison}
\label{sec:quality}
\begin{plan}
  \item \textbf{What each part of the model does} (exact averages, one seed;
        Appendix~\ref{app:model-search}).
        On CelebA $32\times32$: the visible energy alone gives
        \FIDablPixelsOnly{}, one generic face-like texture on every
        trajectory; adding the hidden units without the springs gives
        \FIDablNoSprings; the full model \FIDcelebLowExact. The hidden units
        carry the content, the springs smooth neighbouring pixels: without
        them the same faces appear, with speckle. Replacing the quartic
        on-site term by a harmonic one as well gives \FIDablHarmonic. At a
        glance the three sets of images are hard to tell apart, so say what
        the FID measures here: pixel grain, to which the Inception features
        are very sensitive.
        Label the first number as this paper's visible energy, not as the model
        of \citet{firstpaper2025}, which has many more pair couplings.
  \item \textbf{The global unit carries the class.} With the patch units fixed,
        MNIST improves from \FIDmnistSmallK{} ($K=\Kglobal$) over
        \FIDmnistMidK{} ($K=128$) to \FIDmnistExact{} ($K=\KglobalClasses$),
        and Fashion-MNIST from \FIDfashionSmallK{} over \FIDfashionMidK{} to
        \FIDfashionExact; the share of samples that a classifier assigns to a
        class with confidence rises from 72\% to 83\% on MNIST (real digits:
        98\%). On faces the same change is worth 0.8
        (\FIDcelebLowLargeK{} against \FIDcelebLowExact). Faces are one aligned
        class; ten classes need global templates per class.
        State also: at $K=\Kglobal$ the MNIST strokes are clean but a quarter
        of the samples are not a digit, which the FID alone does not show.
  \item \textbf{Samples} (Figure~\ref{fig:samples}): uncurated, next to real
        images, from the model trained with $S=\Slarge$.
  \item \textbf{Earlier models, in the text only, without a ranking.}
        \citet{wangachour2026} report FIDs of \FIDwaMnist{} (MNIST) and
        \FIDwaFashion{} (Fashion-MNIST), averaged over four seeds, for an
        oscillator network trained adversarially with backpropagation through
        the solver. Their paper states neither the number of samples nor the
        FID implementation or reference set. That alone decides the MNIST
        comparison: in our protocol the FID of real test images is
        \floorMnist{} at 2{,}000 samples and 11.1 at 1{,}000. So quote their
        numbers for orientation and say that they are not comparable with
        ours; the same for \citet{liao2022gaussian} (MNIST 17.5).
        Do not quote their numbers for DTM and NLM (107.8 / 112.8 and
        230.5 / 200.8): these are their own re-runs in a reduced
        configuration (for DTM an $80\times80$ grid, four diffusion steps and
        three binary units per grey pixel), not results of the original
        authors.
        Our statement of quality is the distance to the floor of real images
        at the same sample count (Table~\ref{tab:results}).
        \todo{a like-for-like comparison would need their sample count and
        FID code, or their samples; ask the authors, or leave it at this.}
  \item \textbf{Faces.} Among the models proposed for analog hardware
        \citep{wangachour2026,jelincic2025dtm} none reports faces. Do not write ``the first
        faces'' without that qualifier: a digitally trained Gaussian--Bernoulli
        machine generates colour faces at $32\times32$
        \citep{liao2022gaussian}, without an FID.
        \todo{a digital reference on the same grey data, if run.}
  \item \textbf{Nearest training images.} Samples shown above their nearest
        training images (Appendix~\ref{app:more-results}; the panels exist).
        No quantitative memorisation test is run, so do not write ``no
        memorisation'': write that the samples are not copies of their nearest
        training images, and give the capacity argument ($K$ templates against
        the number of training images).
  \item \textbf{What is still wrong.} \todo{one honest sentence: contrast of
        faces, residual malformed digits.}
\end{plan}

% ---------------------------------------------------------------- Figure 3 ---
\begin{figure*}[t]
  \centering
  \placeholderbox{0.235\textwidth}{2.5in}{%
    \textbf{MNIST}\par\smallskip real (left half)\par generated, $S=\Slarge$
    (right half)\par\smallskip uncurated, fixed seed}\hfill
  \placeholderbox{0.235\textwidth}{2.5in}{%
    \textbf{Fashion-MNIST}\par\smallskip real (left half)\par generated,
    $S=\Slarge$ (right half)\par\smallskip uncurated, fixed seed}\hfill
  \placeholderbox{0.49\textwidth}{2.5in}{%
    \textbf{CelebA}\par\smallskip top: $32\times32$, real and generated\par
    bottom: $64\times64$, real and generated\par\smallskip
    \textbf{Figure 3 placeholder.} Optional last row: the same trajectories
    with exact averages, $S=\Slarge$ and $S=\Ssmall$.}
  \caption{\textbf{Generated images on the three data sets,} from the model
    trained with $S=\Slarge$ sampled hidden states. \todo{caption: uncurated;
    which rows are real; protocol in one line.}}
  \label{fig:samples}
\end{figure*}

\subsection{The Steps of the Sampler}
\begin{plan}
  \item \textbf{Quench.} As sampled against quenched
        (Table~\ref{tab:sampler}): the largest single effect.
  \item \textbf{Corrector.} Small at $32\times32$, larger at $64\times64$.
  \item \textbf{Noise is needed.} The probability-flow ODE is worse
        (\FIDablExactOde{} against \FIDcelebLowExact) and sensitive to the scale
        of the force.
  \item \textbf{Time dependence is needed} (control on an earlier model,
        Appendix~\ref{app:model-search}): Langevin dynamics under the final
        energy alone falls into a blank state. Qualitative parallel: a
        higher-order associative memory at finite temperature fails to retrieve
        when started far from every memory \citep{rooke2026stochastic}.
  \item Could move to the appendix if space is short; keep one sentence here.
\end{plan}

% ----------------------------------------------------------------- Table 2 ---
\begin{table}[t]
  \caption{\textbf{Sampler ablation} on CelebA $32\times32$ (FID).}
  \label{tab:sampler}
  \centering
  \small
  \begin{tabular}{lcc}
    \toprule
    Sampler & exact & $S=\Slarge$ \\
    \midrule
    reverse SDE, as sampled       & \FIDablExactRaw       & \FIDablSlargeRaw \\
    + quench                      & \FIDcelebLowExact     & \FIDcelebLowSlarge \\
    + corrector and quench        & \FIDcelebLowExactCorr & \FIDcelebLowSlargeCorr \\
    probability-flow ODE + quench & \FIDablExactOde       & \FIDablSlargeOde \\
    \bottomrule
  \end{tabular}
\end{table}

% =============================================================================
\section{Towards Hardware}
\label{sec:hardware}
\budget{0.6 page; high level, no circuit design}

\begin{plan}
  \item \textbf{A circuit for the bipartite structure exists; lead with it.}
        \citet{bacvanski2025dense} design and simulate an analog dense
        associative memory: visible and hidden neurons joined by one resistive
        crossbar used in both directions, a softmax over the hidden neurons
        built from bipolar transistors, and hidden neurons about ten times
        faster than the visible ones, which gives the same log-sum-exp energy
        on the visible neurons. It is the deterministic limit of one of our
        Potts units with stored weights.
  \item \textbf{What our model needs beyond that circuit.} A noise source for
        the pixels (the bath); couplings that change in time; the quadratic
        couplings $Q$; signed couplings (their conductances are non-negative);
        the learning rule.
  \item \textbf{Each operation is a known analog primitive.}
        Fields: forward read of a conductance array (crossbar)
        \citep{wan2022neurram}.
        Hidden average: an analog softmax element \citep{elfadel1993softmax};
        a sampled state: a winner-take-all circuit with noise.
        The patch level may not need a winner-take-all: independent bistable
        units per patch did almost as well in an earlier model
        (Appendix~\ref{app:model-search}); the rule covers them as the
        two-state case.
        Force on the pixels: transposed read of the same array, so the coupling
        is reciprocal by construction.
        Learning: rank-one updates per clamped image; with a sampled state only
        the winner's row changes.
  \item \textbf{Most of the model is local.} The patch units are wired to
        their patch only; the global unit has \Kglobal{} states for faces and
        \KglobalClasses{} for the ten-class data.
  \item \textbf{The visible units can be linear and uncoupled; say this
        here.} Without any coupling between visible units, and with only a
        harmonic on-site term, the model generates the same faces on the same
        trajectories, with more pixel grain (FID \FIDablHarmonic{} against
        \FIDcelebLowExact; Section~\ref{sec:quality}). All nonlinearity then
        sits in the hidden units and the device is purely bipartite, as the
        circuit of \citet{bacvanski2025dense}. The nearest-neighbour springs
        are an optional smoothing term worth 5.6 FID.
        \todo{the runs share the seed and the initial templates; say so
        wherever ``the same faces'' is used.}
  \item \textbf{Time dependence.} The couplings are stored at \nSlices{} noise
        levels. Do not claim that a few coupling sets suffice. Replacing the
        \nSlices{} sets after training by 5 to 9 sets per group costs 2.5 FID
        on CelebA $32\times32$ for the model trained from samples (none on an
        earlier model trained with exact averages), and it fails on MNIST and
        Fashion-MNIST (FID 234 and about 90), where the visible and global
        parameters change more in time. One sentence, as an open point;
        numbers in Appendix~\ref{app:model-search}.
  \item \textbf{Not shown here.} Explicit dynamics of the hidden units at a
        finite switching rate; the displacement measurements with read noise;
        the precision of the couplings. \todo{one sentence each.}
\end{plan}

% =============================================================================
\section{Related Work}
\label{sec:related}
\budget{0.8 page; this section cannot be cut}

\begin{plan}
  \item \textbf{Generative models for analog and thermodynamic hardware.}
        Denoising thermodynamic models: a chain of Boltzmann machines with
        latent nodes, trained with a clamped and a free phase
        \citep{jelincic2025dtm}.
        Oscillator networks with hidden units and four weight sets, trained
        adversarially with backpropagation through the solver
        \citep{wangachour2026}.
        Langevin computers trained on the likelihood of reversed trajectories
        \citep{whitelam2025generative}, with an online local rule and hidden
        units on a few MNIST prototypes \citep{wangdeng2026online}.
  \item \textbf{Equilibrium propagation for generative models.}
        \citet{scellier2017equilibrium}; denoising score matching with
        symmetric nudges on synthetic data \citep{de2026symmetric}; energy
        matching at CIFAR scale in simulation \citep{gower2026potentialep};
        flow matching on small digit images \citep{gower2026floweqprop}.
        Difference: these compare a free and a nudged equilibrium of the whole
        system, and the gradient is exact for vanishing nudge. Our rule is one
        phase at full clamp and exact there; the response comes from the
        fluctuations of thermal hidden units; and the score is the system's
        own force.
        \todo{check in the full texts: the statement on the nudge limit for
        each of the three, and whether any of them reads fluctuations.}
  \item \textbf{Energy-based models with hidden units.} Contrastive divergence
        \citep{hinton2002training}; exponential family harmoniums
        \citep{welling2004exponential}, of which our model is a member with
        visible--visible springs; score matching for such models
        \citep{hyvarinen2005estimation,hyvarinen2007connections,swersky2011autoencoders,vincent2011connection}.
        That score matching needs no samples from the model, and that its
        gradient for hidden units is a function of their posterior moments at
        data, is known there as mathematics. New here: these moments as
        fluctuation measurements on thermal hidden units, local to each
        coupler, and training from $S$ reads.
  \item \textbf{Say it ourselves: without visible couplings the model is a
        restricted Boltzmann machine with Gaussian visible units.} Each joint
        hidden state then gives an axis-aligned Gaussian, and the hidden units
        are categorical with a state-dependent stiffness.
        \citet{liao2022gaussian} train a single Gaussian--Bernoulli machine of
        this family: 4096 to 10{,}000 binary hidden units, each coupled to all
        pixels, one learned variance per pixel, one set of couplings. It
        generates MNIST from noise at FID 17.5 (sample count not stated) and,
        without a reported FID, Fashion-MNIST and colour CelebA at
        $32\times32$. Its learning rule is the Boltzmann rule: the product of
        the two activities at a coupler with data clamped, minus the same
        product while the machine runs freely. It is local, and per coupler
        simpler than ours, which needs a sum over the unit's pixels and the
        clamp force. But every update needs a negative phase: a chain of 100
        sampling steps with 10 Langevin steps each, started from noise, that
        is, one generation run per training example and update, and the
        update is a difference of two sampled correlations.
        The trade-off to state: they keep one energy and pay with a
        free-running generation per update; we need no negative phase and pay
        with couplings that change in time (\nSlices{} sets). Also different:
        Potts units, patch units with local wiring, a state-dependent
        stiffness, and the continuous-time protocol. Their model is ahead in
        two respects, say so: one coupling set, and colour.
        \todo{decide whether MNIST 17.5 goes into Table~\ref{tab:results} as a
        digitally trained relative.}
  \item \textbf{Dense associative memories and diffusion models.} The score of
        a diffusion model is the gradient of a dense-associative-memory energy
        \citep{hoover2023memory,ambrogioni2024search}; as the training set
        grows, such models pass from retrieving stored images, through
        spurious states, to generalisation \citep{pham2026emergence}. These
        works read trained neural diffusion models as memories of the training
        data. We build the energy explicitly, with few learned templates, and
        show that templates held at data prototypes are far worse than learned
        ones (Appendix~\ref{app:model-search}).
        Dense associative memories with learned memories exist for
        classification \citep{krotov2016dense}; to our knowledge not as a
        diffusion model that generates images.
        In machine-learning terms a learned memory energy is a mixture model,
        and we say so: Gaussian-mixture scores trained on the denoising
        objective are analysed in theory \citep{shah2023learning}; products
        of Gaussian-mixture experts on filter responses are trained over all
        noise levels for image restoration \citep{zach2024product}; mixtures
        on overlapping patches are a classical image prior
        \citep{zoran2011learning}. Mixing training-set patches locally
        predicts what convolutional diffusion models generate
        \citep{kamb2025analytic}; our patch units learn \Mpatch{} templates
        each instead of storing patches.
        Analog circuits for dense associative memories
        \citep{bacvanski2025dense} and the thermodynamic cost of retrieval at
        finite temperature \citep{rooke2026stochastic} treat stored memories
        and retrieval; we treat learning and generation.
  \item \textbf{Diffusion models.} \citet{song2019generative,ho2020denoising};
        the final denoising step and the stochastic sampler are known practice,
        here as physical operations.
  \item \textbf{Physical learning.}
        \citet{stern2021supervised,dillavou2026imperfection}.
  \item \todo{read the three closest preprints in full before writing this
        section; check for newer ones.}
\end{plan}

% =============================================================================
\section{Limitations and Conclusion}
\label{sec:conclusion}
\budget{0.5 page}

\begin{plan}
  \item \textbf{Limitations.} A simulation in the limit of fast hidden units;
        clamp forces and projections exact; grey images up to $64\times64$;
        \paramsCelebLow{} couplings per slice; face contrast too low; the
        comparison with earlier models is across different sizes and unknown
        sample counts.
  \item \textbf{Conclusion.} Three sentences: what the model is, that it trains
        from sampled statistics without loss, what it reaches.
  \item \textbf{Outlook.} A circuit-level study; explicit hidden dynamics;
        colour; the work and entropy production of the generation protocol,
        for which the framework of \citet{rooke2026stochastic} applies.
\end{plan}

